% Chapter 36: a commit message's anatomy under MiniDAQ's commit policy: a subject, a blank line,
% a body wrapped at 72 that says why, and a trailer block that tools read.
\begin{tikzpicture}[x=1cm,y=1cm]
  \node[draw=d2drule,line width=0.8pt,fill=white,font=\ttfamily\scriptsize,align=left,anchor=north west,inner sep=6pt] (msg) at (0,0) {%
    Book interior: Chapter 24 figures\\[3pt]
    Five figures: the FIFO as a timing diagram, Gray code across a\\
    clock domain, the FIFO through a burst, a run's edges, and how\\
    long a stop takes, from measurements/ch24/run\_sequence.txt.\\[3pt]
    Review: change control checklist, items 1 to 6.\\
    Co-Authored-By: Claude Opus 5.5 <noreply@anthropic.com>};
  \draw[partcolor,line width=1.2pt] (-0.15,-0.12) -- (-0.15,-0.5); \node[note,text=partcolor,anchor=east,align=right] at (-0.25,-0.3) {subject: what,\\at most 72};
  \draw[d2dcat0,line width=1.2pt] (-0.15,-0.65) -- (-0.15,-1.55); \node[note,text=d2dcat0,anchor=east,align=right] at (-0.25,-1.1) {body: why and what\\changed, wrapped at 72};
  \draw[d2dcat2,line width=1.2pt] (-0.15,-1.7) -- (-0.15,-2.25); \node[note,text=d2dcat2,anchor=east,align=right] at (-0.25,-1.97) {trailers: read\\by tools};
\end{tikzpicture}
